\documentclass[10pt,a4paper]{amsart}
\usepackage{amsmath,amssymb,amsthm,graphicx}
\usepackage{verbatim}
\newenvironment{reference}{\comment}{\endcomment}
%\newenvironment{reference}{}{}
\newenvironment{slogan}{\comment}{\endcomment}
\newenvironment{history}{\comment}{\endcomment}

% For commutative diagrams we use Xy-pic
\usepackage[all]{xy}

% We use 2cell for 2-commutative diagrams.
\xyoption{2cell}
% Theorem environments.
%
\theoremstyle{plain}
\newtheorem{theorem}[subsection]{Theorem}
\newtheorem{proposition}[subsection]{Proposition}
\newtheorem{lemma}[subsection]{Lemma}

\theoremstyle{definition}
\newtheorem{definition}[subsection]{Definition}
\newtheorem{example}[subsection]{Example}
\newtheorem{exercise}[subsection]{Exercise}
\newtheorem{situation}[subsection]{Situation}

\theoremstyle{remark}
\newtheorem{remark}[subsection]{Remark}
\newtheorem{remarks}[subsection]{Remarks}

\numberwithin{equation}{subsection}

% Macros
%
\def\lim{\mathop{\mathrm{lim}}\nolimits}
\def\colim{\mathop{\mathrm{colim}}\nolimits}
\def\Spec{\mathop{\mathrm{Spec}}}
\def\Hom{\mathop{\mathrm{Hom}}\nolimits}
\def\Ext{\mathop{\mathrm{Ext}}\nolimits}
\def\SheafHom{\mathop{\mathcal{H}\!\mathit{om}}\nolimits}
\def\SheafExt{\mathop{\mathcal{E}\!\mathit{xt}}\nolimits}
\def\Sch{\mathit{Sch}}
\def\Mor{\mathop{\mathrm{Mor}}\nolimits}
\def\Ob{\mathop{\mathrm{Ob}}\nolimits}
\def\Sh{\mathop{\mathit{Sh}}\nolimits}
\def\NL{\mathop{N\!L}\nolimits}
\def\CH{\mathop{\mathrm{CH}}\nolimits}
\def\proetale{{pro\text{-}\acute{e}tale}}
\def\etale{{\acute{e}tale}}
\def\QCoh{\mathit{QCoh}}
\def\Ker{\mathop{\mathrm{Ker}}}
\def\Im{\mathop{\mathrm{Im}}}
\def\Coker{\mathop{\mathrm{Coker}}}
\def\Coim{\mathop{\mathrm{Coim}}}

% Boxtimes
%
\DeclareMathSymbol{\boxtimes}{\mathbin}{AMSa}{"02}

%
% Macros for moduli stacks/spaces
%
\def\QCohstack{\mathcal{QC}\!\mathit{oh}}
\def\Cohstack{\mathcal{C}\!\mathit{oh}}
\def\Spacesstack{\mathcal{S}\!\mathit{paces}}
\def\Quotfunctor{\mathrm{Quot}}
\def\Hilbfunctor{\mathrm{Hilb}}
\def\Curvesstack{\mathcal{C}\!\mathit{urves}}
\def\Polarizedstack{\mathcal{P}\!\mathit{olarized}}
\def\Complexesstack{\mathcal{C}\!\mathit{omplexes}}
% \Pic is the operator that assigns to X its picard group, usage \Pic(X)
% \Picardstack_{X/B} denotes the Picard stack of X over B
% \Picardfunctor_{X/B} denotes the Picard functor of X over B
\def\Pic{\mathop{\mathrm{Pic}}\nolimits}
\def\Picardstack{\mathcal{P}\!\mathit{ic}}
\def\Picardfunctor{\mathrm{Pic}}
\def\Deformationcategory{\mathcal{D}\!\mathit{ef}}
\usepackage{cleveref}
\crefname{originalsection}{original section}{original sections}
\Crefname{originalsection}{Original section}{Original sections}
\makeatletter
\newlabel{section-algebras}{{073Y}{source}{}{sourcepointer-1}{}}
\newlabel{section-algebras@cref}{{[originalsection][0][]073Y}{[1][0][]source}{}{}{}}
\makeatother
\crefname{originalproposition}{original proposition}{original propositions}
\Crefname{originalproposition}{Original proposition}{Original propositions}
\makeatletter
\newlabel{proposition-separable-splitting-field}{{0752}{source}{}{sourcepointer-2}{}}
\newlabel{proposition-separable-splitting-field@cref}{{[originalproposition][0][]0752}{[1][0][]source}{}{}{}}
\makeatother
\crefname{originallemma}{original lemma}{original lemmas}
\Crefname{originallemma}{Original lemma}{Original lemmas}
\makeatletter
\newlabel{lemma-matrix-algebras}{{074D}{source}{}{sourcepointer-3}{}}
\newlabel{lemma-matrix-algebras@cref}{{[originallemma][0][]074D}{[1][0][]source}{}{}{}}
\makeatother
\crefname{originaltheorem}{original theorem}{original theorems}
\Crefname{originaltheorem}{Original theorem}{Original theorems}
\makeatletter
\newlabel{theorem-wedderburn}{{0747}{source}{}{sourcepointer-4}{}}
\newlabel{theorem-wedderburn@cref}{{[originaltheorem][0][]0747}{[1][0][]source}{}{}{}}
\makeatother
\title{Characterizations of finite central simple algebras}
\author{The Stacks Project Authors}
\date{stacks-project:brauer.tex:local-snapshot-20261012; GNU FDL 1.2 or later}
\begin{document}
\maketitle
\noindent\textit{Editorial note.} All delivered mathematical text is the author's own, copied from the pinned Stacks Project chapter file sites-cohomology.tex; no statement or proof step was rewritten, reordered or removed.\par\medskip
\noindent\textit{Editorial note.} The proof is the author's complete argument; each of the six conditions is accounted for in the circle of implications, the only steps attributed elsewhere being named results of the same chapter carried on this card as navigation stubs.\par\medskip
\noindent\textit{Editorial note.} Every labelled reference occurring inside the counted spans is carried in assembly.external\_references as a non-load-bearing navigation pointer (same-chapter ones at their definition site, cross-chapter ones at the invocation inside the counted span), so no load-bearing mathematics is replaced by a stub. No bibliography citation occurs inside the counted spans.\par\medskip
\noindent\textit{Editorial note.} The same-name check for this card scanned every staged Stacks chapter of this repository -- all sources/stacks-*/pinned-original-source.tex files, the card's own chapter excluded -- and the labels found are recorded in dedup.same\_name\_cross\_chapter on the card.\par\medskip
\noindent\textit{Editorial note.} The preamble is the chapter's own retained minimal preamble (sources/stacks-sites-cohomology/retained-preamble.tex), a verbatim excerpt of the native preamble: the Xy-pic setup and the native theorem declarations with the mathematical macros.\par\medskip
\noindent\textit{Allowed standard prerequisites (editorial).}\begin{itemize}
\item Finite central simple algebras. Same chapter, Section section-algebras and the definitions of a finite central simple algebra. Supplies the equivalence of the first two conditions as a consequence of the definitions.
\item Existence of a separable splitting field. Same chapter, Proposition proposition-separable-splitting-field. Supplies that a finite central simple algebra has a separable splitting field, whose Galois closure is then also a splitting field, giving the chain to the matrix algebra conditions.
\item Wedderburn's theorem. Same chapter, Theorem theorem-wedderburn. Supplies the description of a finite simple algebra as a matrix algebra over a finite skew field, which is the last condition.
\item The matrix algebra case. Same chapter, Lemma lemma-matrix-algebras. Supplies that the base change of a matrix algebra is again central simple, which is used to return from the third condition to the first.
\end{itemize}\medskip
\ifdefined\hypertarget\hypertarget{sourcepointer-1}{}\fi
\noindent\textit{Source reference.} Original section 073Y: sources/stacks-brauer/pinned-original-source.tex, source lines 247--247. Navigation pointer invoked at proof line 1006; the counted span contains the pointer only and the cited result is not reproduced inside it, so the pointer is registered as a non-load-bearing navigation stub in the established form for Stacks material. No bibliography citation is involved.\par\medskip
\ifdefined\hypertarget\hypertarget{sourcepointer-2}{}\fi
\noindent\textit{Source reference.} Original proposition 0752: sources/stacks-brauer/pinned-original-source.tex, source lines 923--923. Navigation pointer invoked at proof line 1008; the counted span contains the pointer only and the cited result is not reproduced inside it, so the pointer is registered as a non-load-bearing navigation stub in the established form for Stacks material. No bibliography citation is involved.\par\medskip
\ifdefined\hypertarget\hypertarget{sourcepointer-3}{}\fi
\noindent\textit{Source reference.} Original lemma 074D: sources/stacks-brauer/pinned-original-source.tex, source lines 477--477. Navigation pointer invoked at proof line 1014; the counted span contains the pointer only and the cited result is not reproduced inside it, so the pointer is registered as a non-load-bearing navigation stub in the established form for Stacks material. No bibliography citation is involved.\par\medskip
\ifdefined\hypertarget\hypertarget{sourcepointer-4}{}\fi
\noindent\textit{Source reference.} Original theorem 0747: sources/stacks-brauer/pinned-original-source.tex, source lines 358--358. Navigation pointer invoked at proof line 1017; the counted span contains the pointer only and the cited result is not reproduced inside it, so the pointer is registered as a non-load-bearing navigation stub in the established form for Stacks material. No bibliography citation is involved.\par\medskip
\begin{lemma}
\label{lemma-finite-central-simple-algebra}
Let $k$ be a field. For a $k$-algebra $A$ the following are equivalent
\begin{enumerate}
\item $A$ is finite central simple $k$-algebra,
\item $A$ is a finite dimensional $k$-vector space, $k$ is the center of $A$,
and $A$ has no nontrivial two-sided ideal,
\item there exists $d \geq 1$ such that
$A \otimes_k \bar k \cong \text{Mat}(d \times d, \bar k)$,
\item there exists $d \geq 1$ such that
$A \otimes_k k^{sep} \cong \text{Mat}(d \times d, k^{sep})$,
\item there exist $d \geq 1$ and a finite Galois extension $k'/k$
such that
$A \otimes_k k' \cong \text{Mat}(d \times d, k')$,
\item there exist $n \geq 1$ and a finite central skew field $K$
over $k$ such that $A \cong \text{Mat}(n \times n, K)$.
\end{enumerate}
The integer $d$ is called the {\it degree} of $A$.
\end{lemma}

\begin{proof}
The equivalence of (1) and (2) is a consequence of the definitions, see
Section \ref{section-algebras}.
Assume (1). By
Proposition \ref{proposition-separable-splitting-field}
there exists a separable splitting field $k \subset k'$ for $A$.
Of course, then a Galois closure of $k'/k$ is a splitting field also.
Thus we see that (1) implies (5). It is clear that (5) $\Rightarrow$ (4)
$\Rightarrow$ (3). Assume (3). Then $A \otimes_k \overline{k}$
is a finite central simple $\overline{k}$-algebra for example by
Lemma \ref{lemma-matrix-algebras}.
This trivially implies that $A$ is a finite central simple $k$-algebra.
Finally, the equivalence of (1) and (6) is Wedderburn's theorem, see
Theorem \ref{theorem-wedderburn}.
\end{proof}

\end{document}
